\section*{Source Code}

\subsection*{Development Environment}
Development and testing are conducted on a dedicated virtual machine (\texttt{dirtyuvm}) running:
\begin{itemize}
    \item Linux kernel version 6.8 (custom build)
    \item Modified NVIDIA UVM driver version 580.65.06 (loaded as a kernel module)
\end{itemize}

\subsection*{Repository}
The full source code is hosted at:
\begin{center}
    \url{https://github.com/aleph-5/open-gpu-kernel-modules}
\end{center}
This is a fork of NVIDIA's open-source GPU kernel modules, with our modifications applied on top.

\subsection*{Modified Driver Files}
All driver modifications reside under \texttt{kernel-open/nvidia-uvm/}. The files containing functionality added or significantly modified by us are:

\begin{enumerate}
    \item \texttt{uvm\_common.h} / \texttt{uvm\_common.c} - The primary location of our implementation. Defines \texttt{struct dirty\_page\_info} (page number, timestamp, PID) and \texttt{struct uvm\_dirty\_page\_table} (an \texttt{xarray}-based per-process table). Implements the core API: table initialisation and teardown, fault recording, and page lookup. Also declares \texttt{uvm\_dirty\_procfs\_init} for the \texttt{procfs} interface.
    \item \texttt{uvm\_va\_block.c} - Hooked to invoke the dirty tracking record function when a write fault is serviced on a VA block.
    \item \texttt{uvm\_gpu\_replayable\_faults.c} - Modified to record dirty page events when GPU write faults are replayed.
    \item \texttt{uvm\_va\_space.c} / \texttt{uvm\_va\_space.h} - Extended to register the dirty-page GPU-unmap invalidation callback (\texttt{uvm\_va\_space\_dirty\_init}) and manage per-process tracking table lifetime alongside VA space creation and destruction.
    \item \texttt{uvm.c} - Top-level integration point; initialises the procfs dirty tracking interface at module load.
\end{enumerate}

\subsection*{Tests}
User-space tests are categorized into correctness and performance tests and located in the \texttt{correctness\_tests/} and \texttt{performance\_tests/} directories respectively

The correctness tests are run everytime the driver is modified and reloaded using  \texttt{run\_regression.sh}